\scanpage{22}
\uline{Corollaire} 2 Soit $K$ extension de type fini de $\mathbb R$,
$X$ un modèle. Conditions équivalentes (i) (i bis) (ii) (ii bis)
(ii ter) (ii quater) (iii) (iii bis).
\deleted{(i) $-1$ une somme de carrés dans $K$}

\uline{Corollaire} 3 $X$ préschéma de type fini sur $\mathbb R$.

Conditions équivalentes

(i) $X(\mathbb R)=\varnothing$

(ii) $\mathrm{cd}_2(X)<+\infty$

\deleted{\unsure{Alors}} Lorsqu'il en est ainsi, alors pour \inserted{tt} faisceau
(étale) $F$ sur $X$, on a
\[
H^i(X,F)\simeq H^i(\underbrace{X(C)/\mathbb Z/2\mathbb Z}_{\text{topologie ordinaire}},F)
\]
et on a
\[
\mathrm{cd}_2(X)\leq2\dim X
\]

\uline{Théorème} 2 Soit $k$ un corps,\inserted{\figure{\begin{tikzpicture}[baseline=(n.base)]\node[draw,ellipse,inner sep=3pt] (n) {$r$ un \deleted{\unsure{entier}} \inserted{\unsure{nb} premier}};\end{tikzpicture}}}
\deleted{$k'$ une \unsure{e}}\inserted{Conditions} équivalentes

\margin{Conditions th. 2 \unsure{satisfaites} \unsure{pourvu que} \uline{l'une}
des conditions suiv. \unsure{ait lieu}
$\left\{\begin{array}{l}
\operatorname{car}k>0\\
r\ne2\\
-1\text{ somme de carrés}\\
\text{\unsure{non}}\ \mathrm{cd}_2k(\sqrt{-1})<+\infty]\\\relax
[\mathrm{cd}_r(k)<+\infty]
\end{array}\right.$}

\figure{\begin{tikzpicture}[baseline=(n.base)]\node[draw,circle,inner sep=2pt,line width=1pt] (n) {i};\end{tikzpicture}}
Pour \unsure{tte|tt} extension de type fini $k'$ de $k$, on a
\[
\mathrm{cd}_r(k')=\mathrm{cd}_r(k)+\operatorname{deg\,tr}k'/k
\]

(i bis) Comme dessus, \deleted{$k'$ une extension quadratique}\inserted{avec $k'=k(\sqrt{-1})$ !}
\deleted{(et même $k'=k(\sqrt{-1})$ !)}

(ii) $k'$ contient un corps\inserted{$k_0$} de type fini sur le corps
premier, et tel que $\mathrm{cd}_r(k_0)<+\infty$, ou $\mathrm{cd}_2k(\sqrt{-1})=+\infty$

\figure{\begin{tikzpicture}[baseline=(n.base)]\node[draw,circle,inner sep=2pt,line width=1pt] (n) {iii};\end{tikzpicture}}
On n'est pas dans le cas suivant: $k$ est
ordonnable (\deleted{donc}\inserted{i.e.} de car $0$\inserted{et $-1$ n'est pas somme de carrés}), $r=2$,
\figure{\begin{tikzpicture}[baseline=(n.base)]\node[draw,ellipse,inner sep=4pt] (n) {$\mathrm{cd}_2(k(\sqrt{-1}))<+\infty$.};\end{tikzpicture}}
\inserted{[d'où $\mathrm{cd}_2(k)=+\infty$]}
\margin{21}

